% Work-MAT | 2. rocnik | pripravovana hodina 22
% Návaznost: uživatelské připravy.zip (hodiny 01–20), autorská učebnice na GitHubu.
% Realisticky.cz: didaktická inspirace v posloupnosti a typech úloh; zadání jsou vlastní.
\documentclass[11pt,a4paper]{article}
\input{../../sablony/pripravy/priprava.sty}
\nastavPripravu{2. ročník}{Exponenciální rovnice}{22}{Kvadratická substituce}
\begin{document}
\blok{Motivace a návaznost}
Zopakuj pomocnou proměnnou $t=2^x>0$ a potom ukaž, že $4^x=(2^x)^2=t^2$. Tím se z exponenciální rovnice stane známá kvadratická.
\poznamka{Orientační nabídka pro 45--55 minut: krátké zopakování a motivace, nový výklad, tři společné vzorové úlohy, samostatná práce. Vyber podle tempa třídy; zbývající zadání slouží jako rezerva.}
\blok{Výklad}
Pomocná rovnice může mít záporný kořen, který nelze přijmout, protože $t=a^x>0$. Oba kroky — výběr přípustných $t$ a návrat k $x$ — zapisujeme odděleně.
\blok{Společné příklady}
\begin{prikladReseniBox}[title={Příklad 1}]
Vyřeš, resp. urči: $4^x-5\cdot2^x+4=0$.
\tcblower
\[\begin{aligned}t=2^x>0,\quad t^2-5t+4=0 \\ (t-1)(t-4)=0\;\Longrightarrow\;t\in\{1,4\} \\ 2^x=1\;\lor\;2^x=4\;\Longrightarrow\;x\in\{0,2\}\end{aligned}\]
\end{prikladReseniBox}
\begin{prikladReseniBox}[title={Příklad 2}]
Vyřeš, resp. urči: $9^x+2\cdot3^x-15=0$.
\tcblower
\[\begin{aligned}t=3^x>0,\quad t^2+2t-15=0 \\ (t-3)(t+5)=0\;\Longrightarrow\;t=3 \\ 3^x=3\;\Longrightarrow\;x=1\end{aligned}\]
\end{prikladReseniBox}
\begin{prikladReseniBox}[title={Příklad 3}]
Vyřeš, resp. urči: $4^x-6\cdot2^x+8=0$.
\tcblower
\[\begin{aligned}t=2^x>0,\quad (t-2)(t-4)=0 \\ 2^x=2\;\lor\;2^x=4 \\ K=\{1,2\}\end{aligned}\]
\end{prikladReseniBox}
\poznamka{Požaduj, aby žáci napsali obě substituce vedle sebe: $t=2^x$ a $4^x=t^2$. Vyřazení záporného kořene musí být výslovné.}
\blok{Samostatné procvičování}
\begin{zadaniReseniBox}[title={Samostatná práce 1}]
\begin{enumerate}[label=\alph*)]
\item $4^x-3\cdot 2^x+2=0$.
\item $9^x-4\cdot3^x+3=0$.
\item $4^x+2^x-2=0$.
\item $25^x-6\cdot5^x+5=0$.
\item $4^x-10\cdot2^x+16=0$.
\item $9^x+3^x-12=0$.
\end{enumerate}
\tcblower
\begin{enumerate}[label=\alph*)]
\item $K=\{0,1\}$.
\item $K=\{0,1\}$.
\item $x=0$.
\item $K=\{0,1\}$.
\item $K=\{1,3\}$.
\item $x=1$.
\end{enumerate}
\end{zadaniReseniBox}
\blok{Dobrovolné rozšíření}
\blok{Úkol na doma}
Při početních úlohách napiš postup, zvolenou metodu a zkoušku v původním zadání. U soustav kontroluj obě rovnice; u nerovnic výsledek zakresli na osu.
\begin{zadaniReseniBox}[title={Domácí úkol}]
\begin{enumerate}[label=\alph*)]
\item $4^x-5\cdot2^x+4=0$.
\item $9^x-10\cdot3^x+9=0$.
\item $25^x-6\cdot5^x+5=0$.
\item $4^x-6\cdot2^x+8=0$.
\item $9^x+2\cdot3^x-15=0$.
\item $4^x-3\cdot2^x+2=0$.
\item Co vyřadíme, vyjde-li $t=-2$?.
\item Kdy se hodí kvadratická substituce?.
\end{enumerate}
\tcblower
\begin{enumerate}[label=\alph*)]
\item $K=\{0,2\}$.
\item $K=\{0,2\}$.
\item $K=\{0,1\}$.
\item $K=\{1,2\}$.
\item $K=\{1\}$.
\item $K=\{0,1\}$.
\item Záporný kořen pomocné rovnice..
\item Při $a^{2x}$ a $a^x$ společného základu..
\end{enumerate}
\end{zadaniReseniBox}
\poznamka{Domácí úkol obsahuje 8 položek; předpokládaná domácí práce přibližně 15--30 minut. Po větších blocích jsou vědomě promíchané dříve probrané typy.}
\end{document}
